\section{Sealed types e pattern matching}
\label{sec:sealed-pattern}

Três JEPs se juntam aqui. O \citetitle{doc:jep409} (final no JDK~17)
deixa o autor de uma classe ou interface restringir quem pode
estendê-la. O \citetitle{doc:jep394} (JDK~16) trouxe o
\texttt{instanceof} com variável de padrão, e o \citetitle{doc:jep441}
(JDK~21) levou os padrões ao \texttt{switch}, que desde o
\citetitle{doc:jep361} (JDK~14) também pode ser expressão.

\subsection*{Um \texttt{switch} exaustivo sobre um tipo selado}

\lstinputlisting[language=Java]{generated/source/SealedForma.java}

O \texttt{switch} de \texttt{area} não tem \texttt{default}, e compila
porque o compilador sabe que \texttt{Forma} só admite \texttt{Circulo} e
\texttt{Quadrado}. Essa informação está no \texttt{.class} da interface,
no atributo \texttt{PermittedSubclasses} (JVMS~§4.7.31,
\citetitle{doc:jvms25}):

\inputbytecode{SealedForma_Forma}

E o método \texttt{area}:

\inputbytecode{SealedForma}

\paragraph{Lendo a listagem.}
\begin{itemize}
\item Offsets 0--5: \texttt{Objects.requireNonNull} no seletor. Sem
  \texttt{case null}, um seletor \texttt{null} lança
  \texttt{NullPointerException}, \enquote{como antes}, nas palavras do
  JEP~441.
\item Offset 11: \texttt{invokedynamic typeSwitch} devolve o índice do
  primeiro \texttt{case} que casa (0 para \texttt{Circulo}, 1 para
  \texttt{Quadrado}), e o \texttt{lookupswitch} salta para ele.
\item Offsets 44--53: o \texttt{default} que \emph{você não escreveu}.
  O \texttt{javac} criou um que faz \texttt{new MatchException} e
  \texttt{athrow}. O JEP~441 explica o motivo: com compilação separada,
  uma nova implementação pode aparecer em tempo de execução, e então o
  \texttt{switch} lança em vez de seguir em silêncio.
\item Depois do salto, \texttt{checkcast} e a chamada ao accessor
  (\texttt{raio()}, \texttt{lado()}): o \enquote{cast seguro} que o padrão
  promete é um \texttt{checkcast} que já sabemos que passa.
\end{itemize}

\paragraph{A decisão que isso muda.}
Quando o conjunto de variantes é fechado e conhecido (resultado de um
pagamento: aprovado, recusado, pendente; comando de uma fila), modele com
\texttt{sealed interface} e records, e trate com \texttt{switch} sem
\texttt{default}. Se alguém acrescentar uma variante, todo
\texttt{switch} que não a trata deixa de compilar: o erro aparece no
build, não em produção. Um \texttt{default} \enquote{por segurança} joga
fora essa garantia.

\paragraph{O bug que isso explica.}
Se a biblioteca que define \texttt{Forma} ganhar \texttt{Triangulo} e o
seu código não for recompilado, o \texttt{default} sintético da listagem
lança \texttt{MatchException}. Não é um bug do \texttt{switch}: é o
sinal de que o binário está desatualizado em relação à hierarquia.

\subsection*{Guardas e ordem dos \texttt{case}}
\lstinputlisting[language=Java]{generated/source/PatternSwitch.java}

\inputbytecode{PatternSwitch}

Repare nos offsets 44--60. O \texttt{case Integer i when i > 10} entra
no índice 0; se a guarda falha (\texttt{if\_icmpgt} não salta), o código
grava \texttt{1} em \texttt{istore\_2} e volta ao offset 9, chamando
\texttt{typeSwitch} de novo a partir do \texttt{case} seguinte. A ordem
dos \texttt{case} é a ordem de teste: um \texttt{case Integer i} antes
do \texttt{case Integer i when \ldots} tornaria o segundo inalcançável,
e o \texttt{javac} recusa isso como erro de dominância.
